% ================================================================
% CHAPTER 1: THE GEOMETRY OF IDEAS
% ================================================================
\chapter{The Geometry of Ideas}
\label{ch:geometry}

\epigraph{Whereof one cannot speak, thereof one must be silent.}{Ludwig Wittgenstein, \textit{Tractatus Logico-Philosophicus}, 7}

\section{Why Embed Ideas in Hilbert Space?}

The central question of this book is deceptively simple: \emph{can the
relationships between ideas be given a geometry?} Not a metaphorical geometry
--- not ``the landscape of ideas'' as a figure of speech --- but a genuine
mathematical geometry, in which ideas are points (or vectors) in a space, their
mutual affinity is measured by angles and distances, and the operations we
perform on ideas (combining them, comparing them, transmitting them) correspond
to well-defined geometric transformations.

The answer we shall develop is: \emph{yes, and the natural geometry is that of
a real Hilbert space.} An idea $a$ is mapped to a vector $\embed{a}$ in a
Hilbert space $\Hspace$, the resonance between ideas becomes an inner product,
and composition of ideas becomes vector addition. This simple pair of axioms
yields a remarkably rich theory --- one in which the Cauchy--Schwarz inequality
limits cross-cultural understanding, persuasion reduces to orthogonal
projection, and Nash equilibria in communication games admit geometric
characterization.

But before we state the axioms, we must explain why this particular mathematical
structure is the right one. This requires a brief tour through the history of
attempts to ``geometrize meaning.''

\subsection{Historical Antecedents}

The idea that meanings have geometric structure is not new. It has appeared in
at least four independent traditions, each approaching the problem from a
different direction.

\paragraph{Osgood's Semantic Differential (1957).}
Charles Osgood, George Suci, and Percy Tannenbaum developed the \emph{semantic
differential}, a technique for measuring the ``meaning'' of a concept by asking
subjects to rate it on bipolar scales (good--bad, strong--weak, active--passive).
Factor analysis of these ratings consistently revealed three principal
dimensions --- evaluation, potency, and activity --- suggesting that meanings
inhabit a low-dimensional space. Osgood's work was empirical rather than
axiomatic: the geometry emerged from data, not from first principles. But it
established the crucial insight that \emph{meaning has dimensionality}, and that
the dimensionality is surprisingly low.

\paragraph{G\"ardenfors's Conceptual Spaces (2000).}
Peter G\"ardenfors proposed that concepts are represented as regions in a
multi-dimensional ``conceptual space,'' with each dimension corresponding to a
quality (color, size, pitch, etc.). The similarity between concepts is inversely
related to their distance in this space. G\"ardenfors's framework is richer than
Osgood's: it accommodates structured concepts (with correlations between
dimensions), prototype effects, and conceptual change. But it remains primarily
a framework for \emph{concepts} (categories applied to external objects) rather
than for \emph{ideas} (propositions, theories, worldviews). IDT~v2 extends
G\"ardenfors's geometric intuition from concepts to ideas.

\paragraph{Latent Semantic Analysis (1990).}
Scott Deerwester and colleagues introduced \emph{Latent Semantic Analysis}
(LSA), which represents words and documents as vectors in a space derived from
the singular value decomposition of a term-document matrix. LSA showed that
purely distributional information --- which words co-occur in which contexts ---
is sufficient to recover semantic relationships: synonyms cluster together,
antonyms are separated, and analogical relationships appear as parallelogram
structures. LSA provided the first large-scale evidence that semantic structure
is genuinely geometric.

\paragraph{Word2Vec and Neural Embeddings (2013).}
Tomas Mikolov and colleagues demonstrated that neural network models trained on
large text corpora produce word embeddings with remarkable algebraic properties:
$\text{king} - \text{man} + \text{woman} \approx \text{queen}$. This ``word
arithmetic'' suggests that the embedding space has genuine linear structure ---
that semantic relationships correspond to vector operations. The explosion of
work on neural embeddings (GloVe, BERT, GPT) has confirmed that distributional
semantics yields embeddings with rich geometric properties. But these embeddings
are learned from data, not derived from axioms. IDT~v2 asks: \emph{what axioms
would make such embeddings principled?}

\subsection{From Empirical to Axiomatic}

Each of these traditions discovered geometric structure in meaning empirically.
Osgood found it in factor analysis, G\"ardenfors postulated it theoretically,
LSA extracted it from co-occurrence matrices, and word2vec learned it from
neural networks. But none provided a set of \emph{axioms} that would explain
\emph{why} meaning has geometric structure, or what constraints this structure
must satisfy.

IDT~v2 fills this gap. We begin with ideas as primitive objects, equipped with
composition and a notion of void, and we postulate that there exists an embedding
into a Hilbert space that preserves these structures. The theorems that follow
are consequences of these axioms --- not of any particular dataset, not of any
particular embedding algorithm, but of the \emph{structural requirements} that
any meaning-preserving embedding must satisfy.

This is the characteristic move of axiomatic mathematics: to identify the
essential properties of a structure and derive consequences that hold for
\emph{all} structures satisfying those properties. Just as group theory studies
all groups (not just permutation groups or matrix groups), IDT~v2 studies all
Hilbert embeddings of ideatic spaces (not just word2vec or GloVe or any
particular implementation).


\section{The Axioms of HilbertIdeaticSpace}

We now state the axioms precisely. Throughout this book, $\IS$ denotes an
ideatic space in the sense of IDT~v1: a set of ideas equipped with a
binary composition operation $\compose$, an identity element $\void$ (the void),
and a depth function $\depth : \IS \to \mathbb{N}$.

\begin{definition}[Hilbert Ideatic Space]
\label{def:hilbert-ideatic-space}
A \textbf{Hilbert Ideatic Space} is a tuple $(\IS, \compose, \void, \depth,
\Hspace, \varphi)$ where:
\begin{enumerate}
\item $(\IS, \compose, \void)$ is a monoid (ideas compose associatively, with
$\void$ as identity);
\item $\Hspace$ is a real Hilbert space;
\item $\varphi : \IS \to \Hspace$ is a function (the \textbf{meaning embedding})
satisfying:
\begin{enumerate}
\item $\embed{\void} = 0$ \hfill (void maps to the zero vector)
\item $\embed{a \compose b} = \embed{a} + \embed{b}$ for all $a, b \in \IS$
\hfill (composition maps to addition)
\end{enumerate}
\item $\depth : \IS \to \mathbb{N}$ is a subadditive function with $\depth(\void)
= 0$.
\end{enumerate}
The \textbf{resonance strength} between ideas $a$ and $b$ is defined as:
\[
  \rs{a}{b} \;:=\; \ip{\embed{a}}{\embed{b}}_\Hspace.
\]
\end{definition}

Let us unpack these axioms one by one, explaining both their mathematical
content and their philosophical motivation.

\begin{axiom}[Monoid Structure]
\label{ax:monoid}
$(\IS, \compose, \void)$ is a monoid: composition is associative, and $\void$
is a two-sided identity.
\end{axiom}

This axiom is inherited from IDT~v1 and requires little additional justification.
Ideas compose: one can combine the idea of ``natural selection'' with the idea
of ``cultural transmission'' to form the idea of ``memetic evolution.'' This
composition is associative (the order of combining three ideas does not matter,
though the order of \emph{two} does). The void $\void$ is the empty idea --- the
idea with no content --- which, when composed with any idea, leaves it unchanged.

\begin{axiom}[Hilbert Space]
\label{ax:hilbert}
$\Hspace$ is a real Hilbert space: a complete inner product space over $\mathbb{R}$.
\end{axiom}

Why a Hilbert space, rather than a mere vector space, a Banach space, or a
metric space? The answer is the inner product. The inner product gives us three
things simultaneously: a notion of \emph{angle} (and hence of similarity and
opposition), a notion of \emph{distance} (and hence of nearness and separation),
and a notion of \emph{projection} (and hence of persuasion and influence). No
weaker structure provides all three.

The completeness requirement (that Cauchy sequences converge) is primarily
technical: it ensures that limits of sequences of meanings exist within the
space, so that we can discuss convergent processes of meaning-change. In
finite-dimensional applications, completeness is automatic.

\begin{axiom}[Void Embedding]
\label{ax:void-embed}
$\embed{\void} = 0$.
\end{axiom}

The void --- the idea with no content --- maps to the zero vector. This is the
unique vector with zero norm (zero ``weight'' or ``significance'') and zero
inner product with every other vector (zero resonance with everything). Silence,
in the Hilbert space of meaning, is the origin.

This axiom has deep philosophical resonance with Wittgenstein's famous dictum:
``Whereof one cannot speak, thereof one must be silent.'' The void is precisely
that about which nothing can be said --- and in the Hilbert space, it says
nothing, resonates with nothing, and occupies no position except the origin.

\begin{axiom}[Composition as Addition]
\label{ax:compose-add}
For all $a, b \in \IS$: $\embed{a \compose b} = \embed{a} + \embed{b}$.
\end{axiom}

This is the most consequential axiom of IDT~v2. It asserts that the embedding
$\varphi$ is a \emph{monoid homomorphism} from $(\IS, \compose, \void)$ to
$(\Hspace, +, 0)$: composition in the ideatic space corresponds to addition in
the Hilbert space.

What does this mean philosophically? It means that ideas \emph{add up} in
meaning space. When you compose two ideas --- say, ``democracy'' and ``market
economics'' --- the resulting composite idea occupies the position in meaning
space that is the vector sum of the two components' positions. The meaning of
the composite is, geometrically, the \emph{sum} of the meanings of the parts.

\begin{interpretation}[Additivity and Compositionality]
\label{interp:additivity}
The composition-as-addition axiom is a strong form of \emph{compositionality}:
the meaning of a complex expression is determined by the meanings of its parts
and the mode of their combination. In IDT~v2, the ``mode of combination'' is
always addition --- the simplest possible combining rule.

This is a deliberate idealization. Real ideas do not always compose additively:
irony, metaphor, and context-dependence all introduce non-linearities. But the
linear approximation is the natural starting point, just as linear approximation
is the natural starting point in physics (Hooke's law, Ohm's law, the harmonic
oscillator). The non-linear refinements belong to the next generation of the
theory.

From the perspective of Frege's principle of compositionality, our axiom is a
particularly clean implementation: the meaning function $\varphi$ is literally
a homomorphism. Frege held that the meaning of a sentence is a function of the
meanings of its parts; we make this precise by requiring that the function be
addition. This is stronger than Frege's principle (which allows arbitrary
functions), but the strength is precisely what gives the theory its power.

The additivity axiom also connects to the tradition of \emph{distributional
semantics}. The word2vec model, for instance, represents phrases approximately
as sums of word vectors. Our axiom asserts that this is not merely an
approximation but a structural requirement: any embedding that respects the
monoid structure of ideas \emph{must} map composition to addition.
\end{interpretation}


\section{The Meaning Embedding as Homomorphism}

The axioms of Definition~\ref{def:hilbert-ideatic-space} together imply that
$\varphi$ is a monoid homomorphism. Let us record this explicitly.

\begin{proposition}[$\varphi$ is a Monoid Homomorphism]
\label{prop:phi-homomorphism}
The embedding $\varphi : (\IS, \compose, \void) \to (\Hspace, +, 0)$ is a
monoid homomorphism. That is:
\begin{enumerate}
\item $\embed{\void} = 0$ (preserves identity);
\item $\embed{a \compose b} = \embed{a} + \embed{b}$ for all $a, b \in \IS$
(preserves the operation).
\end{enumerate}
\end{proposition}

\begin{proof}
Both properties are axioms (Axioms~\ref{ax:void-embed} and~\ref{ax:compose-add}).
The proposition merely observes that these two axioms together constitute the
definition of a monoid homomorphism from $(\IS, \compose, \void)$ to the
additive monoid $(\Hspace, +, 0)$.

In Lean~4, this is formalized as:
\begin{lstlisting}
theorem phi_monoid_hom (a b : Idea) :
    φ (a ∘ b) = φ a + φ b := embed_compose a b
\end{lstlisting}
\end{proof}

A monoid homomorphism preserves the algebraic structure of composition, but it
need not be injective. Two distinct ideas may map to the same vector in
$\Hspace$. This is not a defect but a feature: it allows the theory to model
\emph{synonymy} (distinct ideas with the same meaning) and \emph{paraphrase}
(different ways of expressing the same thought).

\begin{definition}[Synonymy]
\label{def:synonymy}
Two ideas $a, b \in \IS$ are \textbf{synonymous} (in the Hilbert space model)
if $\embed{a} = \embed{b}$. We write $a \equiv_\varphi b$.
\end{definition}

\begin{proposition}[Synonymy is a Congruence]
\label{prop:synonymy-congruence}
The relation $\equiv_\varphi$ is a congruence on $(\IS, \compose)$: if $a
\equiv_\varphi a'$ and $b \equiv_\varphi b'$, then $a \compose b \equiv_\varphi
a' \compose b'$.
\end{proposition}

\begin{proof}
If $\embed{a} = \embed{a'}$ and $\embed{b} = \embed{b'}$, then:
\[
  \embed{a \compose b} = \embed{a} + \embed{b} = \embed{a'} + \embed{b'}
  = \embed{a' \compose b'}.
\]
The first and third equalities use Axiom~\ref{ax:compose-add}; the middle
equality uses the hypotheses.

In Lean~4:
\begin{lstlisting}
theorem synonymy_congruence (ha : φ a = φ a') (hb : φ b = φ b') :
    φ (a ∘ b) = φ (a' ∘ b') := by
  simp [embed_compose, ha, hb]
\end{lstlisting}
\end{proof}

\begin{interpretation}[Synonymy and the Quotient Space]
\label{interp:synonymy}
The congruence $\equiv_\varphi$ allows us to form the quotient monoid $\IS /
{\equiv_\varphi}$, in which synonymous ideas are identified. This quotient
embeds \emph{injectively} into $\Hspace$, and its image is the set of ``realized
meanings'' --- the vectors in $\Hspace$ that actually correspond to ideas.

From a philosophical perspective, the quotient construction formalizes the
distinction between \emph{sense} (Sinn) and \emph{reference} (Bedeutung) in
Frege's philosophy of language. Two expressions may have different senses (they
are distinct elements of $\IS$) but the same reference (they map to the same
vector in $\Hspace$). The Hilbert space $\Hspace$ is the space of references;
the ideatic space $\IS$ is the space of senses. The embedding $\varphi$ is the
function that sends each sense to its reference.

This also connects to Quine's thesis of the indeterminacy of translation. If
the embedding $\varphi$ is not unique --- if there are multiple valid embeddings
of the same ideatic space into Hilbert space --- then the ``meaning'' of an idea
is not fully determined by its structural relationships. Different embeddings
would identify different pairs of ideas as synonymous, leading to different
quotient structures. This is precisely Quine's point: there is no fact of the
matter about whether two expressions in different languages have the ``same
meaning,'' because ``same meaning'' depends on the choice of embedding.
\end{interpretation}


\section{First Consequences of the Axioms}

The axioms already yield several non-trivial consequences. We collect the most
immediate ones here.

\begin{proposition}[Iterated Composition]
\label{prop:iterated-compose}
For any idea $a \in \IS$ and any $n \in \mathbb{N}$, define $a^{\langle n
\rangle}$ (the $n$-fold self-composition of $a$) inductively by $a^{\langle 0
\rangle} = \void$ and $a^{\langle n+1 \rangle} = a \compose a^{\langle n
\rangle}$. Then:
\[
  \embed{a^{\langle n \rangle}} = n \cdot \embed{a}.
\]
\end{proposition}

\begin{proof}
By induction on $n$. The base case $n = 0$: $\embed{a^{\langle 0 \rangle}} =
\embed{\void} = 0 = 0 \cdot \embed{a}$. The inductive step: assuming
$\embed{a^{\langle n \rangle}} = n \cdot \embed{a}$,
\begin{align*}
  \embed{a^{\langle n+1 \rangle}}
  &= \embed{a \compose a^{\langle n \rangle}} \\
  &= \embed{a} + \embed{a^{\langle n \rangle}} \\
  &= \embed{a} + n \cdot \embed{a} \\
  &= (n+1) \cdot \embed{a}.
\end{align*}

In Lean~4:
\begin{lstlisting}
theorem embed_composeN (a : Idea) (n : ℕ) :
    φ (composeN a n) = n • φ a := by
  induction n with
  | zero => simp [composeN, embed_void, zero_smul]
  | succ n ih => simp [composeN, embed_compose, ih, succ_nsmul]
\end{lstlisting}
\end{proof}

This result has a beautiful geometric interpretation: repeating an idea $n$
times moves you $n$ times as far along the same direction in meaning space. The
meaning of ``justice, justice, justice'' is three times the vector of ``justice''
--- it has three times the norm (three times the emphasis, three times the
weight) but the same direction (the same qualitative meaning).

\begin{proposition}[Resonance of Iterated Ideas]
\label{prop:resonance-iterated}
For any ideas $a, b \in \IS$ and $n, m \in \mathbb{N}$:
\[
  \rs{a^{\langle n \rangle}}{b^{\langle m \rangle}} = nm \cdot \rs{a}{b}.
\]
\end{proposition}

\begin{proof}
Direct computation:
\begin{align*}
  \rs{a^{\langle n \rangle}}{b^{\langle m \rangle}}
  &= \ip{n \cdot \embed{a}}{m \cdot \embed{b}} \\
  &= nm \cdot \ip{\embed{a}}{\embed{b}} \\
  &= nm \cdot \rs{a}{b}.
\end{align*}
Here we used Proposition~\ref{prop:iterated-compose} and the bilinearity of the
inner product.
\end{proof}

\begin{corollary}[Self-Resonance of Iterated Ideas]
\label{cor:self-resonance-iterated}
$\rs{a^{\langle n \rangle}}{a^{\langle n \rangle}} = n^2 \cdot \rs{a}{a}$.
In other words, the ``weight'' of an idea grows quadratically with repetition.
\end{corollary}

\begin{proof}
Set $b = a$ and $m = n$ in Proposition~\ref{prop:resonance-iterated}.
\end{proof}

\begin{interpretation}[Repetition and Emphasis]
\label{interp:repetition}
The quadratic growth of self-resonance under repetition has a striking
interpretation in rhetoric and propaganda. When a message is repeated $n$
times, its ``weight'' (self-resonance) does not merely double or triple ---
it grows as $n^2$. This formalizes the intuition behind political slogans,
advertising jingles, and religious mantras: repetition does not merely
\emph{add} emphasis, it \emph{multiplies} it.

This quadratic growth is a consequence of the linear embedding: if
$\embed{a^{\langle n \rangle}} = n \cdot \embed{a}$, then the norm squared
is $n^2 \nrm{\embed{a}}^2$. The non-linearity (quadratic growth of weight
from linear growth of position) is characteristic of inner product spaces
and distinguishes the Hilbert space model from simpler models (such as
counting repetitions, which would give linear growth).

In Bourdieu's terms, this suggests that cultural capital accumulates
quadratically: an agent who repeats their message $n$ times acquires not
$n$ but $n^2$ units of cultural weight. This accords with the observation
that dominant cultural discourses are not merely more frequent but
\emph{disproportionately} more influential than their frequency alone would
predict.
\end{interpretation}

\begin{proposition}[Composition and Resonance]
\label{prop:compose-resonance}
For any ideas $a, b, c \in \IS$:
\[
  \rs{a \compose b}{c} = \rs{a}{c} + \rs{b}{c}.
\]
\end{proposition}

\begin{proof}
\begin{align*}
  \rs{a \compose b}{c}
  &= \ip{\embed{a \compose b}}{\embed{c}} \\
  &= \ip{\embed{a} + \embed{b}}{\embed{c}} \\
  &= \ip{\embed{a}}{\embed{c}} + \ip{\embed{b}}{\embed{c}} \\
  &= \rs{a}{c} + \rs{b}{c},
\end{align*}
using linearity of the inner product in the first argument.

In Lean~4:
\begin{lstlisting}
theorem rs_compose_left (a b c : Idea) :
    rs (a ∘ b) c = rs a c + rs b c := by
  simp [rs, embed_compose, inner_add_left]
\end{lstlisting}
\end{proof}

\begin{interpretation}[Linearity of Cultural Evaluation]
\label{interp:linearity}
Proposition~\ref{prop:compose-resonance} states that the resonance of a
composite idea with any third idea is the \emph{sum} of the resonances of the
parts. This is a strong claim about cultural evaluation: when a culture
evaluates a composite idea (say, ``democratic socialism''), its evaluation
is the sum of its evaluations of the parts (``democracy'' and ``socialism''
separately).

This linearity assumption is again an idealization. In practice, the
combination of two ideas can produce emergent effects that are not captured by
simple addition --- irony, paradox, and dialectical synthesis are all phenomena
where the whole is more (or less, or different) than the sum of the parts. But
the linear approximation provides a baseline against which these non-linear
effects can be measured, just as linear regression provides a baseline against
which non-linear relationships can be detected.

The proposition also implies that resonance is a \emph{linear functional} of the
first argument (for fixed second argument). In functional-analytic terms, each
idea $c$ defines a continuous linear functional $\rs{\cdot}{c}$ on the image of
$\varphi$. By the Riesz representation theorem, this functional is represented
by the vector $\embed{c}$ --- which is precisely the definition we started with.
This circularity is reassuring: it means the definition is \emph{consistent}
with the ambient Hilbert space structure.
\end{interpretation}


\section{The Image of the Embedding}

Not every vector in $\Hspace$ need be the image of an idea. The image
$\varphi(\IS) = \{\embed{a} : a \in \IS\}$ is a subset of $\Hspace$ that
inherits structure from both the ideatic space and the Hilbert space.

\begin{proposition}[Image is an Additive Submonoid]
\label{prop:image-submonoid}
The image $\varphi(\IS)$ is an additive submonoid of $\Hspace$: it contains
$0$ and is closed under addition.
\end{proposition}

\begin{proof}
$0 = \embed{\void} \in \varphi(\IS)$. If $\embed{a}, \embed{b} \in
\varphi(\IS)$, then $\embed{a} + \embed{b} = \embed{a \compose b} \in
\varphi(\IS)$.
\end{proof}

The image need not be a subspace (it need not be closed under scalar
multiplication) nor a closed subset (limits of images need not be images). These
are important structural features:

\begin{remark}[The Image Need Not Be a Subspace]
\label{rmk:image-not-subspace}
Since $\IS$ is a monoid, not a group, there is in general no idea $a^{-1}$ such
that $a \compose a^{-1} = \void$. Correspondingly, $\varphi(\IS)$ need not be
closed under negation: $-\embed{a}$ need not lie in the image. This means that
for every idea $a$, there is a ``direction of meaning'' (namely $-\embed{a}$)
that corresponds to no idea at all --- a meaning that is, as it were,
\emph{unsayable}.

Moreover, the image is not closed under scalar multiplication in general. The
vector $\frac{1}{2}\embed{a}$ --- ``half'' of idea $a$ --- may correspond to no
idea. You cannot take half a thought.
\end{remark}

\begin{interpretation}[Unsayable Directions]
\label{interp:unsayable}
The existence of directions in $\Hspace$ that are not in the image of $\varphi$
gives precise mathematical content to the notion of the \emph{unsayable}.
Wittgenstein distinguished between what can be \emph{said} (expressible in
propositions) and what can only be \emph{shown} (manifest in the structure of
language but not statable within it). In IDT~v2, the ``said'' corresponds to
$\varphi(\IS)$ --- the set of realizable meanings --- while the ``shown''
corresponds to the ambient structure of $\Hspace$ that is not captured by any
particular idea.

The negative direction $-\embed{a}$ is a particularly suggestive case. It
represents the ``negation'' of $a$ in meaning space --- a direction of meaning
that is maximally opposed to $a$. If this direction is not in the image of
$\varphi$, then the negation of $a$ is \emph{unsayable}: one can say $a$, but
there is no single idea that represents ``not $a$'' in the same geometric
sense. This accords with the observation that negation in natural language is
often parasitic on affirmation: ``not guilty'' derives its meaning from
``guilty,'' and there is no single word that captures the ``opposite'' of
a complex idea like ``Enlightenment rationalism.''

This connects also to Derrida's notion of \emph{diff\'erance}: meaning is
constituted by differences within a system, not by positive terms. The Hilbert
space $\Hspace$ is the system of differences; the image $\varphi(\IS)$ is the
set of positive terms. The structure of meaning is determined as much by what
\emph{cannot} be said (the complement of $\varphi(\IS)$ in $\Hspace$) as by
what can.
\end{interpretation}


\section{The Void as Origin}

The void plays a distinguished role in the Hilbert space model. Let us collect
its properties.

\begin{proposition}[Properties of the Void]
\label{prop:void-properties}
The void $\void$ satisfies:
\begin{enumerate}
\item $\embed{\void} = 0$ (the zero vector);
\item $\nrm{\embed{\void}} = 0$ (zero norm);
\item $\rs{\void}{a} = 0$ for all $a \in \IS$ (zero resonance with everything);
\item $d(\void, a) = \nrm{\embed{a}}$ for all $a \in \IS$ (distance from void
equals norm).
\end{enumerate}
\end{proposition}

\begin{proof}
(1) is Axiom~\ref{ax:void-embed}. For (2), $\nrm{0} = 0$ by definition of
the norm. For (3), $\rs{\void}{a} = \ip{0}{\embed{a}} = 0$ by linearity. For
(4), $d(\void, a) = \nrm{\embed{\void} - \embed{a}} = \nrm{0 - \embed{a}} =
\nrm{\embed{a}}$.

In Lean~4:
\begin{lstlisting}
theorem rs_void_left (a : Idea) : rs void a = 0 := by
  simp [rs, embed_void, inner_zero_left]
\end{lstlisting}
\end{proof}

\begin{interpretation}[Silence and the Origin]
\label{interp:silence}
The void is the origin of meaning space. Every other idea $a$ is located at
distance $\nrm{\embed{a}}$ from the void --- and this distance is precisely
the idea's norm, its ``weight'' or ``significance.'' An idea with large norm
is far from silence; an idea with small norm is close to it.

This gives precise content to the intuition that some ideas are ``weightier''
than others. The Communist Manifesto, Darwinian evolution, and the theory of
general relativity are ideas of enormous norm --- they are far from the origin,
they resonate strongly with themselves ($\rs{a}{a} = \nrm{\embed{a}}^2$), and
they project significantly onto many other cultural vectors. A shopping list,
by contrast, has small norm: it is close to the origin, weakly self-resonant,
and projects insignificantly onto the great ideas of civilization.

But norm is not the same as \emph{depth} (combinatorial complexity), as we
shall see in Chapter~\ref{ch:depth-norm}. A haiku may have small depth (few
compositional layers) but enormous norm (great semantic weight). A legal
contract may have great depth (many nested clauses and conditions) but modest
norm (limited cultural significance). The independence of these two measures is
one of the key insights of the Hilbert space model.
\end{interpretation}


\section{Wittgenstein's Picture Theory and the Hilbert Space}

We conclude this chapter with a philosophical discussion that frames the
entire project.

In the \emph{Tractatus Logico-Philosophicus}, Wittgenstein proposed that
propositions are \emph{pictures} of reality: a proposition represents a possible
state of affairs by sharing its ``logical form.'' The relationship between a
proposition and what it represents is not one of resemblance but of
\emph{structural correspondence}: the elements of the proposition are correlated
with elements of reality, and the arrangement of elements in the proposition
mirrors the arrangement of elements in reality.

IDT~v2's Hilbert space embedding can be understood as a mathematical
implementation of the picture theory. Ideas (propositions, theories, worldviews)
are mapped to vectors in $\Hspace$, and these vectors ``picture'' the aspects
of reality that the ideas are about. The inner product measures the degree of
structural correspondence between two pictures: two ideas that picture the same
aspect of reality (or closely related aspects) have a large positive inner
product; two ideas that picture unrelated aspects have zero inner product; two
ideas that picture contradictory aspects have a negative inner product.

The composition axiom ($\embed{a \compose b} = \embed{a} + \embed{b}$) then
says that combining two pictures corresponds to adding their vectors --- a
natural operation if pictures are understood as representing \emph{features}
of reality that can be superimposed.

\begin{interpretation}[From Pictures to Vectors]
\label{interp:wittgenstein}
Wittgenstein's picture theory was revolutionary but ultimately unsatisfying:
it told us \emph{that} propositions picture reality, but not \emph{how} to
compute with pictures, measure their similarity, or analyze their composition.
The Hilbert space model fills these gaps. Vectors are pictures we can compute
with: we can measure their angle, compute their sum, project one onto another,
and decompose them into orthogonal components.

The passage from pictures to vectors is, in a sense, the passage from
philosophy to mathematics. Wittgenstein saw that meaning has structure;
IDT~v2 provides the tools to analyze that structure quantitatively.

But Wittgenstein himself might have objected to this mathematization. In the
\emph{Philosophical Investigations}, he abandoned the picture theory in favor
of a view of meaning as \emph{use}: the meaning of a word is its use in a
language game, not its correspondence to a picture or a vector. The tension
between the ``picture'' (vector) view and the ``use'' (game-theoretic) view
is, in fact, productive: Part~II of this book develops the game theory of
meaning, showing how the geometric structure of Part~I interacts with
strategic considerations of communication and persuasion. The two views are
not contradictory but complementary, just as wave mechanics and matrix
mechanics turned out to be equivalent formulations of quantum theory.
\end{interpretation}

\begin{example}[Three Ideas in the Plane]
\label{ex:three-ideas}
To build intuition, consider three ideas embedded in $\Hspace = \mathbb{R}^2$:
\begin{itemize}
\item $a$ = ``democracy,'' with $\embed{a} = (3, 1)$;
\item $b$ = ``free markets,'' with $\embed{b} = (2, 2)$;
\item $c$ = ``authoritarianism,'' with $\embed{c} = (-2, 1)$.
\end{itemize}
Then:
\begin{align*}
  \rs{a}{b} &= \ip{(3,1)}{(2,2)} = 6 + 2 = 8 > 0 \quad \text{(positive
  resonance --- aligned)}, \\
  \rs{a}{c} &= \ip{(3,1)}{(-2,1)} = -6 + 1 = -5 < 0 \quad \text{(negative
  resonance --- opposed)}, \\
  \rs{b}{c} &= \ip{(2,2)}{(-2,1)} = -4 + 2 = -2 < 0 \quad \text{(mildly
  opposed)}.
\end{align*}
The composite idea $a \compose b$ = ``democratic free markets'' has embedding
$(3,1) + (2,2) = (5,3)$, with norm $\sqrt{34} \approx 5.83$ --- greater than
either component alone.

The composition $a \compose c$ = ``democratic authoritarianism'' (an oxymoron,
perhaps, but a genuine political concept in some analyses) has embedding
$(3,1) + (-2,1) = (1,2)$, with norm $\sqrt{5} \approx 2.24$ --- \emph{less}
than either component. The opposition between the components causes partial
cancellation, reducing the weight of the composite idea. This is a geometric
model of \emph{cognitive dissonance}: when you try to hold two opposing ideas
simultaneously, the result has less weight than either alone.
\end{example}

\begin{example}[The Embedding Need Not Be Injective]
\label{ex:non-injective}
Consider two ideas: $a$ = ``the morning star'' and $b$ = ``the evening star.''
These are distinct ideas (different senses, in Frege's terminology), but they
refer to the same object (Venus). In a Hilbert embedding, they might satisfy
$\embed{a} = \embed{b}$ --- they are synonymous in the embedding, even though
they are distinct as ideas. The embedding captures reference (Bedeutung), not
sense (Sinn).

Alternatively, an embedding might assign them nearby but distinct vectors,
capturing the fact that they have slightly different connotations even though
they refer to the same object. The choice depends on the particular embedding
and the granularity of meaning one wishes to capture.
\end{example}

\bigskip

This chapter has laid the foundations: ideas embed in a Hilbert space, composition
maps to addition, and the void maps to the origin. In the next chapter, we
develop the theory of resonance --- the inner product between ideas --- and
discover its deep connections to semiotics, cultural affinity, and the limits of
understanding.
